\begin{lemma}[Exact compact polytope of source-private kernels]\label{lem:step-004-kernel-polytope}
Under the primitive finite one-block definitions, let \(d_*>0\), fix
integers \(N\ge2\), \(M\ge8\), and put
\(\bar\delta_M=d_*/(M^2\log M)\). Define

\[
\mathcal S_{N,M}:=([N]\times\{0,1\})^M,
\qquad
\mathcal G_N:=\{0,1\}^{[N]},
\qquad
D:=|\mathcal S_{N,M}|\,|\mathcal G_N|=(2N)^M2^N.
\]

Let \(\mathcal K_{N,M}\subset\mathbb R^D\) consist of the arrays
\(K=(K_{s,g})_{s\in\mathcal S_{N,M},g\in\mathcal G_N}\) satisfying

\[
K_{s,g}\ge0,
\qquad
\sum_{g\in\mathcal G_N}K_{s,g}=1
\quad(s\in\mathcal S_{N,M}),
\tag{6}
\]

and, for every \emph{ordered} adjacent pair \(s\simeq s'\) and every event
\(E\subseteq\mathcal G_N\),

\[
\sum_{g\in E}K_{s,g}
\le
e^{0.1}\sum_{g\in E}K_{s',g}+\bar\delta_M.
\tag{7}
\]

Then \(\mathcal K_{N,M}\) is a nonempty compact convex polytope, and its
points are in exact bijection with the output kernels of randomized
unrestricted \((0.1,\bar\delta_M)\)-DP maps
\(B:\mathcal S_{N,M}\to\mathcal G_N\).
\end{lemma}

\begin{proof}
Both \(\mathcal S_{N,M}\) and \(\mathcal G_N\) are finite. Consequently,
(6) is a finite family of linear equalities and inequalities. Although the
all-event privacy family is large, it is also finite: there are finitely many
ordered adjacent pairs and exactly \(2^{|\mathcal G_N|}\) subsets of the
finite output space. Thus (7) is the \emph{exact} all-event approximate-DP
definition expressed as finitely many closed affine halfspaces. In
particular, no invalid replacement of all-event privacy by singleton-only
constraints is made.

The product of row simplexes defined by (6) is closed and bounded in
\(\mathbb R^D\), hence compact. Intersecting it with the finitely many closed
halfspaces (7) leaves a compact set. All constraints are affine, so the
intersection is convex and is a polytope.

For nonemptiness, fix any \(g_0\in\mathcal G_N\) and set

\[
K^0_{s,g}:=\mathbf 1\{g=g_0\}
\quad\text{for every }s,g.
\]

The output distribution is identical for \(s\) and \(s'\). Hence it is
\((0,0)\)-DP and satisfies (7), including the events containing \(g_0\) and
those not containing it. Thus \(K^0\in\mathcal K_{N,M}\).

If \(B\) is any randomized map, define
\(K_{s,g}=\Pr[B(s)=g]\). Row-stochasticity is (6), and its all-event DP
inequality is exactly (7), so a source-private \(B\) gives a point of
\(\mathcal K_{N,M}\). Conversely, given \(K\in\mathcal K_{N,M}\), on input
\(s\) draw \(g\) from the finite row distribution \((K_{s,g})_g\). This
defines a randomized unrestricted map \(B_K\), and (7) proves its exact
source-cap privacy for every event. The full set \(\mathcal G_N\) includes
every improper, nonmonotone hypothesis.

Only the output law on each input matters to both privacy and the one-call
expected loss. Therefore arbitrary internal randomization is encoded without
loss. Also, if a public object is hardwired into a learner's code, the
resulting map still has one such finite kernel and belongs to
\(\mathcal K_{N,M}\) whenever it satisfies source privacy. If a map is
\((\varepsilon,\delta)\)-DP with
\(0\le\varepsilon\le0.1\) and \(0\le\delta\le\bar\delta_M\), its event
inequality is at most the right side of (7), so its kernel also belongs to
\(\mathcal K_{N,M}\). Thus the polytope includes the exact cap and every
stronger privacy pair below it. \end{proof}

\begin{lemma}[Continuous payoff geometry and uniform strict game value]\label{lem:step-004-payoff-value}
Under Lemma~\ref{lem:step-004-kernel-polytope} and Proposition~\ref{prop:step-003-expected-hardness}, fix integers
\(N\ge N_*\), \(M\ge8\) satisfying (H). Give \([N+1]\) the discrete
topology and \(\Delta([N])\) its Euclidean simplex topology, and set
\(\mathcal I_N=[N+1]\times\Delta([N])\). Then:

\begin{enumerate}
\item \(\mathcal I_N\) is compact.
\item The expected one-block loss \(\ell_{N,M}(K;t,Q)\) admits the finite
   coefficient representation
   \[
   \ell_{N,M}(K;t,Q)=\langle K,a_{N,M}(t,Q)\rangle.
   \]
   It is affine and
   continuous in \(K\), and jointly continuous in \((K,t,Q)\). For fixed
   \(K,t\), its dependence on \(Q\) is generally polynomial of degree at
   most \(M+1\), so only continuity, not affinity in \((t,Q)\), is claimed.
\item The function
   \[
   f_{N,M}(K):=\max_{(t,Q)\in\mathcal I_N}
      \ell_{N,M}(K;t,Q)
   \]
   is continuous on \(\mathcal K_{N,M}\), and the attained value
   \[
   v_{N,M}:=\min_{K\in\mathcal K_{N,M}}f_{N,M}(K)
   \tag{8}
   \]
   satisfies \(v_{N,M}>\eta\).
\end{enumerate}
\end{lemma}

\begin{proof}
The simplex \(\Delta([N])\) is closed and bounded in \(\mathbb R^N\), and
\([N+1]\) is finite. Their product \(\mathcal I_N\) is compact. This includes
every boundary face and every point mass of the simplex.

Write an arbitrary labeled dataset as

\[
s=((x_1,y_1),\ldots,(x_M,y_M))\in\mathcal S_{N,M}.
\]

For \((t,Q)\in\mathcal I_N\) and \(g\in\mathcal G_N\), define

\[
p_{t,Q}(s)
:=\left(\prod_{j=1}^M Q(x_j)\right)
  \mathbf 1\{y_j=\tau_t(x_j)\text{ for every }j\},
\tag{9}
\]

and

\[
r_{t,Q}(g)
:=R_Q(g,\tau_t)
=\sum_{x=1}^N Q(x)\mathbf 1\{g(x)\ne\tau_t(x)\}.
\tag{10}
\]

Thus \(p_{t,Q}\) is exactly the mass function of
\((Q^{\tau_t})^M\), including zero mass on every label-inconsistent
dataset. Define the coefficient vector \(a_{N,M}(t,Q)\in\mathbb R^D\) by

\[
a_{N,M}(t,Q)_{s,g}:=p_{t,Q}(s)r_{t,Q}(g).
\tag{11}
\]

For the learner represented by \(K\), direct expansion of the sample and
output expectations gives

\[
\begin{aligned}
\ell_{N,M}(K;t,Q)
&:=\mathbb E_{\substack{S\sim(Q^{\tau_t})^M\\g\sim B_K(S)}}
       R_Q(g,\tau_t)\\
&=\sum_{s\in\mathcal S_{N,M}}p_{t,Q}(s)
  \sum_{g\in\mathcal G_N}K_{s,g}r_{t,Q}(g)\\
&=\sum_{s,g}K_{s,g}a_{N,M}(t,Q)_{s,g}
=\langle K,a_{N,M}(t,Q)\rangle.
\end{aligned}
\tag{12}
\]

For fixed \(t\), every coordinate in (11) is a polynomial in the coordinates
of \(Q\): (9) has degree \(M\) and (10) has degree one. Hence
\(a_{N,M}\) is continuous on each simplex copy and therefore on the finite
disjoint union \(\mathcal I_N\). Equation (12) is linear, hence affine and
continuous, in \(K\), and is jointly continuous in \((K,t,Q)\). The product
structure in (9) is precisely why no affine dependence on raw \(Q\) is
asserted or needed.

The compact coefficient image
\(a_{N,M}(\mathcal I_N)\) is bounded. Therefore, for any
\(K,K'\in\mathcal K_{N,M}\),

\[
\begin{aligned}
|f_{N,M}(K)-f_{N,M}(K')|
&\le
\max_{(t,Q)\in\mathcal I_N}
  |\langle K-K',a_{N,M}(t,Q)\rangle|\\
&\le \|K-K'\|_2
\max_{(t,Q)\in\mathcal I_N}\|a_{N,M}(t,Q)\|_2.
\end{aligned}
\tag{13}
\]

Thus \(f_{N,M}\) is continuous. Its inner maximum is attained because
\(\mathcal I_N\) is compact, and its minimum in (8) is attained because
\(\mathcal K_{N,M}\) is compact.

Now fix any \(K\in\mathcal K_{N,M}\). By
Lemma~\ref{lem:step-004-kernel-polytope}, \(K\) defines an arbitrary
randomized unrestricted source-private map \(B_K\).  Proposition
\ref{prop:step-003-expected-hardness} therefore supplies some
\((t,Q)\in\mathcal I_N\) such that
\[
\ell_{N,M}(K;t,Q)>\eta.
\]
Hence

\[
f_{N,M}(K)>\eta
\quad\text{for every }K\in\mathcal K_{N,M}.
\tag{14}
\]

Let \(K_*\) attain the minimum in (8). Applying (14) to this actual
minimizer, rather than taking an unattained infimum, yields

\[
v_{N,M}=f_{N,M}(K_*)>\eta.
\tag{15}
\]

This is the required compactness upgrade from learner-by-learner strict
hardness to one uniform strict game value. \end{proof}

\begin{lemma}[Compact bilinear minimax via finite-dimensional separation]\label{lem:step-004-compact-minimax}
Let \(X\) and \(Y\) be nonempty compact convex subsets of finite-dimensional
real vector spaces, and let \(p:X\times Y\to\mathbb R\) be continuous and
affine in each argument (in particular, a restricted bilinear pairing).
Then both outer extrema below are attained and

\[
\min_{x\in X}\max_{y\in Y}p(x,y)
=
\max_{y\in Y}\min_{x\in X}p(x,y).
\tag{16}
\]
\end{lemma}

\begin{proof}
Continuity on the compact product implies that
\(x\mapsto\max_{y\in Y}p(x,y)\) is continuous; for example, uniform
continuity of \(p\) bounds the change of the maximum by the same uniform
modulus. Hence

\[
\alpha:=\min_{x\in X}\max_{y\in Y}p(x,y)
\]

is attained. The elementary weak minimax inequality gives

\[
\max_{y\in Y}\min_{x\in X}p(x,y)\le\alpha.
\tag{17}
\]

We prove the reverse inequality by separation. Fix any real \(r<\alpha\).
For each \(x\in X\), compactness of \(Y\) supplies some \(y_x\in Y\) with
\(p(x,y_x)>r\). Continuity in \(x\) gives a neighborhood on which this same
strict inequality holds. Compactness of \(X\) yields a finite subcover, so
there are \(y_1,\ldots,y_m\in Y\) such that

\[
\max_{1\le j\le m}p(x,y_j)>r
\quad\text{for every }x\in X.
\tag{18}
\]

Define the affine continuous map

\[
F:X\to\mathbb R^m,
\qquad
F(x)=(p(x,y_1),\ldots,p(x,y_m)),
\]

and the closed convex lower orthant

\[
\mathcal D_r:=(-\infty,r]^m.
\]

The image \(F(X)\) is compact and convex, and (18) says
\(F(X)\cap\mathcal D_r=\varnothing\). Strong separation provides a nonzero
\(\lambda\in\mathbb R^m\) oriented so that

\[
\sup_{z\in\mathcal D_r}\langle\lambda,z\rangle
<
\inf_{x\in X}\langle\lambda,F(x)\rangle.
\tag{19}
\]

Every coordinate of \(\lambda\) must be nonnegative: if
\(\lambda_j<0\), sending the \(j\)-th coordinate of \(z\in\mathcal D_r\)
to \(-\infty\) makes the left side of (19) infinite. Normalize the nonzero
nonnegative vector so that \(\sum_j\lambda_j=1\). Then the left side of
(19) is exactly \(r\), attained at \((r,\ldots,r)\), and hence

\[
\inf_{x\in X}\sum_{j=1}^m\lambda_jp(x,y_j)>r.
\tag{20}
\]

Convexity of \(Y\) puts \(y_r:=\sum_j\lambda_jy_j\) in \(Y\), while
affinity in the second argument turns (20) into

\[
\min_{x\in X}p(x,y_r)>r.
\tag{21}
\]

Choose \(r_q\uparrow\alpha\). Compactness of \(Y\) gives a convergent
subsequence \(y_{r_q}\to y_*\in Y\). For every fixed \(x\in X\), (21) and
continuity give

\[
p(x,y_*)=\lim_q p(x,y_{r_q})\ge\lim_q r_q=\alpha.
\]

Thus \(\min_xp(x,y_*)\ge\alpha\). Together with (17), this proves (16) and
shows that the right-hand maximum is attained at \(y_*\). \end{proof}

\begin{lemma}[Compact coefficient hull and exact finite barycenters]\label{lem:step-004-finite-barycenter}
Let \(I\) be a nonempty compact space and let \(a:I\to\mathbb R^D\) be
continuous. Then

\[
\mathcal C:=\operatorname{conv}(a(I))
\]

is compact and convex. Moreover, every \(c\in\mathcal C\) has an exact
representation

\[
c=\sum_{j=1}^r\lambda_j a(i_j),
\qquad
1\le r\le D+1,
\quad
\lambda_j\ge0,
\quad
\sum_{j=1}^r\lambda_j=1,
\tag{22}
\]

with \(i_j\in I\). Consequently, for every \(K\in\mathbb R^D\),

\[
\langle K,c\rangle
=\sum_{j=1}^r\lambda_j\langle K,a(i_j)\rangle
\tag{23}
\]

exactly.
\end{lemma}

\begin{proof}
Convexity is immediate from the definition. For the support bound, start
from any finite convex representation
\(c=\sum_{j=1}^q\lambda_ja(i_j)\), which exists by the definition of convex
hull, and first remove every zero-weight term. Thus all remaining
\(\lambda_j>0\). If \(q>D+1\), the points
\(a(i_1),\ldots,a(i_q)\) are affinely dependent. Hence there are scalars
\(\theta_j\), not all zero, such that

\[
\sum_{j=1}^q\theta_j=0,
\qquad
\sum_{j=1}^q\theta_ja(i_j)=0.
\tag{24}
\]

After changing all signs if needed, at least one \(\theta_j>0\). Put

\[
\rho:=\min_{j:\theta_j>0}\frac{\lambda_j}{\theta_j}.
\]

Then \(\lambda'_j=\lambda_j-\rho\theta_j\) are nonnegative, sum to one,
represent the same \(c\) by (24), and at least one formerly positive
coefficient becomes zero. Removing zero coefficients and iterating gives
(22). This is the affine-dependence proof of Caratheodory's theorem.

The representation bound also proves compactness. Pad every representation
to exactly \(D+1\) terms using zero weights. Then \(\mathcal C\) is the
image of the compact set

\[
\Delta([D+1])\times I^{D+1}
\]

under the continuous map
\((\lambda,i_1,\ldots,i_{D+1})\mapsto
\sum_j\lambda_ja(i_j)\). Thus \(\mathcal C\) is compact. Equation (23)
follows from bilinearity of the Euclidean pairing and introduces no
approximation. \end{proof}

\begin{proposition}[Finite public hard prior for source-private one-block kernels]\label{prop:step-004-finite-hard-prior}
Let \(b_*,d_*>0\) and \(N_*\ge2\) be the constants supplied by Proposition~\ref{prop:step-003-expected-hardness}. Under that proposition and
Lemmas~\ref{lem:step-004-kernel-polytope}--\ref{lem:step-004-finite-barycenter},
for every integer \(N\ge N_*\) and integer \(M\ge8\) satisfying (H), there
is a probability law \(\mu_{N,M}\) supported on at most

\[
D+1=(2N)^M2^N+1
\tag{25}
\]

pairs in \([N+1]\times\Delta([N])\) such that every randomized unrestricted
\((0.1,d_*/(M^2\log M))\)-DP one-block learner satisfies (HP). The prior is
selected independently of the learner and may be revealed to it.
\end{proposition}

\begin{proof}
Use the coefficient map in (11) and define

\[
\mathcal A_{N,M}:=a_{N,M}(\mathcal I_N),
\qquad
\mathcal C_{N,M}:=\operatorname{conv}(\mathcal A_{N,M}).
\]

Lemma~\ref{lem:step-004-finite-barycenter} makes
\(\mathcal C_{N,M}\) a nonempty compact convex subset of \(\mathbb R^D\).
For every fixed \(K\), linearity of the pairing gives

\[
\max_{c\in\mathcal C_{N,M}}\langle K,c\rangle
=
\max_{a\in\mathcal A_{N,M}}\langle K,a\rangle
=f_{N,M}(K).
\tag{26}
\]

Indeed, a convex combination cannot exceed the largest paired coefficient,
and \(\mathcal A_{N,M}\subseteq\mathcal C_{N,M}\) gives the reverse
inequality.

Apply Lemma~\ref{lem:step-004-compact-minimax} with

\[
X=\mathcal K_{N,M},
\qquad
Y=\mathcal C_{N,M},
\qquad
p(K,c)=\langle K,c\rangle.
\]

The first set is nonempty compact convex by
Lemma~\ref{lem:step-004-kernel-polytope}; the second is nonempty compact
convex by Lemma~\ref{lem:step-004-finite-barycenter}; and the payoff is
continuous bilinear. Thus all hypotheses of the separation-based minimax
theorem are discharged. Equations (8), (15), and (26) yield the exact
quantifier swap

\[
\begin{aligned}
v_{N,M}
&=\min_{K\in\mathcal K_{N,M}}
  \max_{c\in\mathcal C_{N,M}}\langle K,c\rangle\\
&=\max_{c\in\mathcal C_{N,M}}
  \min_{K\in\mathcal K_{N,M}}\langle K,c\rangle
>\eta.
\end{aligned}
\tag{27}
\]

Let \(c_*\in\mathcal C_{N,M}\) attain the maximum in (27). By
Lemma~\ref{lem:step-004-finite-barycenter}, there are
\(r\le D+1\), pairs \((t_j,Q_j)\in\mathcal I_N\), and nonnegative weights
\(\lambda_j\) summing to one such that

\[
c_*=\sum_{j=1}^r\lambda_j a_{N,M}(t_j,Q_j).
\tag{28}
\]

Remove zero weights and merge repeated pairs, and define the finite law

\[
\mu_{N,M}:=\sum_{j=1}^r\lambda_j\,\delta_{(t_j,Q_j)}.
\tag{29}
\]

For every \(K\in\mathcal K_{N,M}\), equations (12), (27), and (28) give

\[
\begin{aligned}
\mathbb E_{(t,Q)\sim\mu_{N,M}}\ell_{N,M}(K;t,Q)
&=\sum_{j=1}^r\lambda_j
  \langle K,a_{N,M}(t_j,Q_j)\rangle\\
&=\langle K,c_*\rangle\\
&\ge\min_{K'\in\mathcal K_{N,M}}\langle K',c_*\rangle\\
&=v_{N,M}>\eta.
\end{aligned}
\tag{30}
\]

The first equality is the exact finite-prior expectation, not an
approximation to a continuum. Lemma~\ref{lem:step-004-kernel-polytope}
identifies every source-private randomized unrestricted \(B\) with one such
\(K\), so (30) is precisely (HP).

The construction of \(\mu_{N,M}\) uses the whole fixed strategy set
\(\mathcal K_{N,M}\), not any selected learner. Once (29) is fixed and made
public, a learner whose program refers to (29) still induces some kernel
\(K\in\mathcal K_{N,M}\) if it is source-private; (30) was proved for every
such kernel after the prior was fixed. Hence prior awareness does not change
the quantifier order.

Finally, no interiority was used. The compact instance space contains
\(t=1,N+1\), every boundary distribution, and every point-mass \(Q\).
Equations (9)--(12) remain polynomial and exact when some \(Q(x)=0\).
Different thresholds that agree on the support of a point mass may give
duplicate coefficient vectors, but duplicates neither invalidate the convex
hull nor change (30). \end{proof}
